\chapter{Conclusion}
\label{chap:conclusion}

\section{Project Overview}

This project investigated whether GAN-generated phishing URLs can improve a
URL-based phishing detector and whether the selected detector can be used in a
working browser extension. The work began with real-only classifiers, continued
with several URL generators and controlled augmentation experiments, and ended
with external evaluation and browser implementation.

The final classification approach, HYB-MS, combines a character-level CNN with
35 lexical URL features. The character branch learns local patterns in the URL
text, while the lexical branch represents measurements such as length, digit
and symbol ratios, subdomain structure, path depth, and query structure. The
model was trained as a three-seed ensemble on the registered multi-source real
training data. A matching model was then trained with an additional 10\%
filtered MS-HardGAN-C phishing URLs. Keeping the architecture, real rows,
evaluation data, and calibration procedure fixed allowed the effect of GAN
augmentation to be measured fairly.

The selected model was exported to ONNX and integrated into a Chromium browser
extension. The extension performs local inference, presents several warning
levels, supports English and Arabic, and allows the user to continue or permit
an exact hostname. Reliability controls reduce repeated counting, and
user-evaluation data remain local unless a participant chooses to export them.

\section{Answers to the Research Questions}

\subsection{RQ1: Performance of real-only classifiers}

The real-only experiments showed that URL phishing can be detected effectively
when the training and test distributions are similar. The CNN and LSTM reached
internal F1-scores above 97\%. However, the external evaluations separated the
models more clearly. Some models retained a very low FPR by classifying few
URLs as phishing, but consequently missed most phishing samples. This showed
that accuracy or FPR alone is not sufficient for model selection.

Under equal calibration, LSTM-MS was the strongest ordinary neural baseline.
HYB-MS real-only achieved a better overall external balance because it combined
learned character patterns with full-URL lexical structure. On the primary
30K low-prevalence evaluation, real-only HYB-MS achieved 79.06\% F1 and
81.07\% recall, compared with 74.93\% F1 and 75.20\% recall for LSTM-MS.

\subsection{RQ2: Quality of generated phishing URLs}

The GAN experiments showed that uniqueness alone does not establish quality.
A generator can produce many different strings that are invalid, structurally
abnormal, repetitive in pattern, or easy to distinguish from real phishing
URLs. Generated pools were therefore evaluated using validity, uniqueness,
real-data overlap, lexical realism, diversity, source discrimination, and
their ability to challenge frozen classifiers.

SeqGAN was the strongest of the initial text generators, while MS-HardGAN-C
provided the final filtered augmentation pool. The accepted pool contained no
exact overlap with the permitted real training URLs. The generator did not
predict future zero-day attacks; instead, it recombined patterns learned from
the available phishing training data to provide additional difficult examples.

A later matched-input experiment retrained SeqGAN on the exact same 84,280
phishing rows used by MS-HardGAN-C. This improved SeqGAN's downstream HYB-MS
result, showing that generator input data mattered. MS-HardGAN-C still achieved
the stronger recall, F1, and FNR balance on the frozen 10K and 30K comparisons.
Its selection is therefore supported by a direct input-controlled comparison,
not only by the earlier studies in which the generators used different
phishing pools.

\subsection{RQ3: Effect of GAN augmentation}

GAN augmentation did not improve every model, ratio, or metric. Random
augmentation and higher synthetic percentages frequently reduced performance
on the unchanged real test data. These results are important because they show
that simply increasing the quantity of generated URLs is not a reliable method
of improving detection.

The clearest benefit came from the controlled 10\% MS-HardGAN-C condition. On
the primary 30K external dataset, HYB-MS+10\% achieved 98.12\% accuracy,
80.95\% precision, 81.60\% recall, 81.27\% F1, 96.82\% ROC-AUC, and 87.27\%
PR-AUC. It also reduced FPR from 1.26\% to 1.01\% and FNR from 18.93\% to
18.40\% relative to real-only HYB-MS. The paired F1 improvement was supported
by the reported confidence interval and McNemar comparison, although the
recall change alone was not statistically significant.

The result answers the research question positively but conditionally:
carefully filtered and limited GAN augmentation can improve the balance of URL
phishing detection under the evaluated conditions. The experiments do not show
that augmentation always improves performance.

The matched-generator experiment also showed why the generator choice matters.
On the balanced external 10K set, matched SeqGAN reached 91.59\% F1 and
14.36\% FNR, while MS-HardGAN-C reached 92.67\% F1 and 12.48\% FNR at the
same 1.36\% FPR. Both generators had received the same real phishing training
identities and both supplied exactly 40,410 synthetic rows to HYB-MS.

\subsection{RQ4: Generalisation to different and imbalanced sources}

The selected 10\% model improved F1 and reduced both FPR and FNR relative to
its real-only control on several frozen low-prevalence external datasets. These
included the PhishTank 10K set, the final labelled 6K set, the post-freeze 10K
set, and the primary 30K set. The final labelled 6K evaluation contained only
5\% phishing and produced 98.23\% accuracy, 84.15\% precision, 79.67\% recall,
81.85\% F1, 0.79\% FPR, and 20.33\% FNR for HYB-MS+10\%.

The external scores were lower than the strongest internal scores. This is
evidence of distribution shift rather than a contradiction: later URLs came
from different sources and contained structures that were less represented in
training. The long-URL test also showed a trade-off, where augmentation found
additional phishing URLs but introduced additional false warnings. Therefore,
the model generalises usefully across the tested sets, but it does not remove
distribution shift.

\subsection{RQ5: Browser implementation}

The selected ensemble was successfully deployed as a local browser-extension
prototype. Preprocessing parity, ONNX numerical parity, navigation handling,
localisation, consent, deduplication, policy, and browser smoke tests were
completed. The extension separates the raw model score from the action shown
to the user, allowing moderate evidence to produce a non-blocking notice while
reserving navigation interruption for the strictest tier.

The browser study also showed why extension policy must be evaluated separately
from classifier quality. Blanket trust for Google Workspace suppressed warnings
for phishing hosted in user-controlled documents. Replacing it with a capped,
non-blocking policy recovered 24 detections on the controlled snapshot and
reduced FNR from 21.33\% to 13.33\% without changing the measured FPR. User
feedback further identified false warnings on login pages, long service URLs,
technical sites, Google services, and AI services. The extension satisfies the
project implementation objective as a research prototype, but it is not
presented as a certified production security product.

\section{Achievement of the Project Objectives}

\begin{table}[htbp]
\centering
\caption{Final status of the objectives defined in Chapter~1.}
\label{tab:conclusion-objectives}
\small
\begin{tabular}{@{}p{0.11\textwidth}p{0.64\textwidth}p{0.14\textwidth}@{}}
\toprule
Objective & Evidence completed in the project & Status \\
\midrule
O1 & Label cleaning, canonicalisation, class counts, length analysis, duplicate removal, and identity-disjoint splits. & Achieved \\
O2 & RF, SVM, CNN, RNN, LSTM, and hybrid real-only classifiers trained and evaluated. & Achieved \\
O3 & URL GAN variants trained only with the permitted phishing training data. & Achieved \\
O4 & Validity, uniqueness, overlap, realism, diversity, source, and challenge analyses completed. & Achieved \\
O5 & Real-only and fixed-ratio augmented conditions compared with controlled rows, seeds, and test sets. & Achieved \\
O6 & Internal, source-held-out, later, low-prevalence, PhishTank, long-URL, and GAN-challenge tests completed. & Achieved \\
O7 & Selected ensemble exported to ONNX and implemented in a locally evaluated browser extension. & Achieved \\
\bottomrule
\end{tabular}
\end{table}

All seven objectives were completed within the research scope. Achievement of
the deployment objective means that a functioning browser prototype was built
and evaluated. It does not mean that the system has completed the security,
privacy, device-compatibility, and maintenance work required for unrestricted
public deployment.

\section{Main Findings}

The main findings of the project are as follows:

\begin{itemize}
\item Internal random-split performance can overestimate performance on URLs
from later or different sources.

\item Combining character learning and lexical structure gave the hybrid a
better external balance than the ordinary baselines under equal calibration.

\item Generated URL quality must be evaluated separately from downstream
classifier performance. High uniqueness does not guarantee realistic or useful
augmentation.

\item Large or unfiltered synthetic additions were often harmful. The best
evidence supported a limited, filtered 10\% augmentation condition.

\item The benefit of the selected GAN augmentation was repeated on several
low-prevalence frozen datasets, but its size varied by dataset.

\item Thresholds control the FPR--FNR trade-off. Lower FNR can be obtained by
accepting more false warnings, but no threshold removes distribution overlap.

\item Browser policy can either expose or hide model detections. Trusted
platforms must not silently override strong model evidence on user-controlled
content.

\item User feedback is valuable for finding usability problems, but formal
accuracy metrics require labelled evaluation data rather than incomplete
self-reported outcomes.
\end{itemize}

\section{Limitations}

The detector uses URL text and lexical structure only. It does not inspect the
page DOM, visual design, certificate chain, DNS history, domain age, network
behaviour, or reputation feeds. A convincing phishing URL can therefore be
missed, while a long or unusual legitimate URL can receive a warning.

The evaluation datasets cannot represent every future campaign, language,
country, hosting platform, or browsing pattern. Labels from public feeds can
also contain delays or errors. Although repeated seeds and frozen comparisons
improve reproducibility, the GAN and classifier findings remain specific to
the implemented models and available data.

The browser user evaluation was useful for interface and false-warning
analysis, but it was not a fully supervised clinical-style study. Not every
visited URL received an independently verified label, devices and browsing
patterns differed, and some participants deliberately tested phishing. Its
counts must therefore remain separate from the formal model metrics.

\section{Final Conclusion}

The project provides a valid experimental answer to its main question. GAN
augmentation can improve phishing URL detection when generated samples are
filtered, limited, and compared fairly with a matching real-only model. The
10\% MS-HardGAN-C augmented hybrid produced the strongest and most consistent
external balance in this study. At the same time, the negative augmentation
results and remaining FNR show that synthetic data is not a replacement for
diverse real data.

The completed browser extension demonstrates that the research model can run
locally and support real-time URL warnings. Its tiered decisions, reviewed
service rules, and user controls make it suitable for graduation-project
demonstration and further supervised testing. Continued data collection and
multi-signal verification would be required before treating it as a
production-ready phishing defence.
